% SPDX-License-Identifier: CC-BY-SA-4.0
% Author: Matthieu Perrin

\begingroup

\SetKwBroadcast{FC}{fc}
\SetKwBroadcast{FLC}{flc}

\begin{exercice}[Canaux avec pertes de messages]

  Lorsque les canaux peuvent perdre des messages, aucune propriété de vivacité
  ne peut être garantie si les pertes sont totalement arbitraires.
  On considère donc deux propriétés de vivacité plus faibles que la fiabilité.

  \begin{description}
  \item[Fiabilité]
    Si un processus $p_i$ envoie un message $m$ à un processus correct $p_j$,
    alors $p_j$ finit par recevoir $m$ de $p_i$.

  \item[FC-fiabilité]
    Pour tout message $m$, si un processus $p_i$ envoie $m$ une infinité de fois
    à un processus correct $p_j$, alors $p_j$ reçoit $m$ de $p_i$ une infinité
    de fois.

  \item[FLC-fiabilité]
    Si un processus $p_i$ envoie une infinité de messages à un processus correct
    $p_j$, alors $p_j$ reçoit une infinité de messages de $p_i$.
  \end{description}

  Un canal vérifiant la propriété \textbf{FC-fiabilité} est appelé
  \emph{fair channel} (\FC), et un canal vérifiant la propriété
  \textbf{FLC-fiabilité} est appelé \emph{fair lossy channel} (\FLC).

  \begin{questions}
  \item Montrer que les propriétés \textbf{FC-fiabilité} et
    \textbf{FLC-fiabilité} sont incomparables :
    \begin{enumerate}
    \item donner une exécution qui vérifie \textbf{FC-fiabilité} mais pas
      \textbf{FLC-fiabilité} ;
    \item donner une exécution qui vérifie \textbf{FLC-fiabilité} mais pas
      \textbf{FC-fiabilité}.
    \end{enumerate}

  \item Montrer que chacun de ces deux types de canaux peut être simulé à
    partir de l'autre.
    Pour chaque direction, proposer une simulation et justifier qu'elle
    vérifie la propriété de fiabilité attendue.
  \end{questions}

\end{exercice}

\endgroup
\endinput
